% !TEX program = xelatex
\documentclass[UTF8,a4paper,zihao=5,fontset=fandol]{ctexart}

% 本模板依据《经济研究》网站 2025-01-03 发布的“稿件体例要求”整理。
% 网页已明确的字体字号要求已尽量固化；页边距和行距网页未明示，以下采用中性设置，可按编辑部反馈微调。
% 建议使用 XeLaTeX 编译。

\usepackage[a4paper,left=2.8cm,right=2.8cm,top=2.8cm,bottom=2.6cm,footskip=1.2cm]{geometry}
\usepackage{fontspec}
\usepackage{amsmath,amssymb,bm}
\usepackage{array,booktabs,longtable,tabularx,threeparttable}
\usepackage{graphicx}
\usepackage{caption}
\usepackage{setspace}
\usepackage{indentfirst}
\usepackage{titlesec}
\usepackage{perpage}

\setmainfont{Times New Roman}
\defaultfontfeatures{Ligatures=TeX}

\MakePerPage{footnote}

\setlength{\parindent}{2em}
\setlength{\parskip}{0pt}
\setstretch{1.45}
\pagestyle{plain}

\makeatletter
\renewcommand\footnotesize{\songti\zihao{6}}
\renewcommand{\footnoterule}{%
  \vspace{0.6\baselineskip}%
  \hrule width 0.18\textwidth
  \vspace{0.3\baselineskip}%
}
\makeatother

\newcommand{\ERJTitleSize}{\zihao{-3}} % 需要三号时改为 \zihao{3}
\newcommand{\ERJAuthorSize}{\zihao{-4}}

\DeclareCaptionFont{erjfigurefont}{\heiti\zihao{-5}}
\DeclareCaptionFont{erjtablefont}{\songti\zihao{5}}
\captionsetup{justification=centering,singlelinecheck=false}
\captionsetup[figure]{font=erjfigurefont,labelsep=space}
\captionsetup[table]{font=erjtablefont,labelsep=space,position=top}

\newcolumntype{L}[1]{>{\raggedright\arraybackslash}p{#1}}
\newcolumntype{C}[1]{>{\centering\arraybackslash}p{#1}}
\newcommand{\ERJTableFont}{%
  \fangsong\zihao{-5}%
  \setlength{\tabcolsep}{2pt}%
  \renewcommand{\arraystretch}{1.16}%
}
\newcommand{\ERJTableNoteFont}{\songti\zihao{6}}

\renewcommand{\thesection}{\chinese{section}}
\renewcommand{\thesubsection}{（\chinese{subsection}）}
\renewcommand{\thesubsubsection}{\arabic{subsubsection}.}
\renewcommand{\theparagraph}{(\arabic{paragraph})}
\setcounter{secnumdepth}{4}

\titleformat{\section}
  {\centering\songti\zihao{-4}}
  {\thesection、}
  {0.5em}
  {}
\titlespacing*{\section}{0pt}{1.4\baselineskip}{0.8\baselineskip}

\titleformat{\subsection}[hang]
  {\songti\zihao{5}}
  {\hspace*{2em}\thesubsection}
  {0.3em}
  {}
\titlespacing*{\subsection}{0pt}{0.9\baselineskip}{0.3\baselineskip}

\titleformat{\subsubsection}[hang]
  {\songti\zihao{5}}
  {\hspace*{2em}\thesubsubsection}
  {0.3em}
  {}
\titlespacing*{\subsubsection}{0pt}{0.7\baselineskip}{0.2\baselineskip}

\titleformat{\paragraph}[runin]
  {\songti\zihao{5}}
  {\hspace*{2em}\theparagraph}
  {0.3em}
  {}
\titlespacing*{\paragraph}{0pt}{0.5\baselineskip}{0.5em}

\newcommand{\erjthanks}[1]{\thanks{\hspace*{2em}#1}}
\newcommand{\erjfootnote}[1]{\footnote{\hspace*{2em}#1}}

\newcommand{\ERJAbstract}[1]{%
  \par\noindent{\fangsong\zihao{5}内容摘要：#1\par}%
}

\newcommand{\ERJKeywords}[1]{%
  \par\noindent{\fangsong\zihao{5}关键词：#1\par}%
}

\newcommand{\ERJWhere}[1]{%
  \par\noindent\hspace*{2em}{\songti\zihao{5}其中，#1\par}%
}

\newcommand{\coefse}[2]{%
  \begin{tabular}{@{}c@{}}#1\\[-0.1em] (#2)\end{tabular}%
}

\newcommand{\ERJRef}[1]{%
  \par\noindent\hspace*{2em}\hangindent=2em\hangafter=1 {#1}\par%
}

\makeatletter
\renewcommand{\maketitle}{%
  \begingroup
  \renewcommand{\thefootnote}{\fnsymbol{footnote}}%
  \begin{center}
    {\songti\ERJTitleSize \@title\par}
    \vspace{1.1em}
    {\songti\ERJAuthorSize \@author\par}
  \end{center}
  \vspace{1.0em}
  \@thanks
  \endgroup
  \setcounter{footnote}{0}%
  \renewcommand{\thefootnote}{\arabic{footnote}}%
}
\makeatother

\title{地方债务压力约束下政府产业引导基金与城市创新：实证检验}
\author{%
  作者甲\erjthanks{作者单位，邮政编码：100000，电子信箱：author1@example.com。本文系某基金项目“项目名称”（项目编号：XXXXXX）的阶段性成果。}%
  \quad
  作者乙\erjthanks{作者单位，邮政编码：100000，电子信箱：author2@example.com。}%
}
\date{}

\begin{document}

\maketitle

\ERJAbstract{本文基于2015—2024年地级市—年份面板结果，考察政府产业引导基金能否促进城市创新，以及地方债务压力是否削弱这一促进效应。基准固定效应估计显示，基金设立数量与累计设立规模均与城市专利申请产出显著正相关，其中累计设立规模的结果最为稳定。进一步引入债务压力后，基金累计规模与滞后一期债务压力的交互项在发明申请、实用新型申请和专利申请总量三个口径下均显著为负，说明债务压力会明显削弱基金促进创新的边际作用。机制检验进一步表明，债务压力可能通过三条路径抑制基金绩效：其一，降低基金投向早期项目的比例，削弱耐心资本属性；其二，在低金融发展地区降低社会资本撬动效率；其三，削弱基金改善存款口径金融发展水平、缓解融资约束的能力。需要说明的是，耐心资本机制依赖于基金投资事件数不少于2的分母稳定化样本，社会资本机制主要出现在低金融发展样本中，融资约束机制属于地区金融发展层面的反向代理证据。}
\ERJKeywords{政府产业引导基金；债务压力；城市创新；耐心资本；社会资本撬动效率}

\section{引言}

这里填写整篇论文的引言部分。建议在本节补充研究背景、现实问题、研究空白与本文边际贡献。当前主稿已经接入完整的实证分析章节，后续可在不改变版式的前提下继续扩写。

\section{文献综述与理论假说}

这里填写整篇论文的文献综述与理论假说部分。建议围绕“政府产业引导基金促进创新”“债务压力削弱基金绩效”以及耐心资本、社会资本撬动效率和融资约束缓解三条机制提出可检验假说。

% 本文件由 main.tex 输入，作为“模型设定与变量说明”章节，不单独编译。

\section{模型设定与变量说明}

\subsection{数据来源与样本说明}

本文以地级市—年份为研究单元，样本期为2015—2024年。选择地级市层面展开分析，主要是因为政府产业引导基金通常以注册地、出资地和服务地方产业为重要运行边界，地方债务压力和财政约束也主要在地方政府层面体现。与企业层面数据相比，地级市面板更适合识别地方财政债务约束下政府产业引导基金与区域创新产出之间的关系。

本文使用的创新产出数据来自 CNRDS 地方专利数据，并整理为地级市—年份层面的发明申请量、实用新型申请量和专利申请总量。政府产业引导基金基础信息来自清科基金目录和投中基金目录，本文将两类目录合并去重后，依据基金注册地手工匹配到地级市层面，进而构造各城市基金当年设立数量、累计设立数量和累计设立规模。基金出资结构主要来自投中基金目录中的 LP 和出资人信息，用以识别政府资本、国资平台、政府产业引导基金、FOF 以及其他社会资本出资。基金投资事件信息主要来自清科投资事件数据，本文据此识别政府引导基金的投资时间、投资阶段、投资次数和投资金额，并按基金注册地汇总到地级市—年份层面。地方债务数据依据各地方统计局和统计年鉴披露信息手工整理，控制变量主要来自 CEIC、CNKI 等城市统计资料。

样本处理遵循现有回归面板的统一口径。回归前按城市和年份进行去重，模型中加入城市固定效应和年份固定效应，并按城市维度聚类稳健标准误。由于部分债务变量、基金投资事件变量和机制变量存在缺失，不同回归表中的样本量会随变量口径变化，本文在实证结果表中分别报告各规格观测值数量。

\subsection{变量定义}

本文被解释变量为城市创新产出。正文主结果使用专利申请量而非授权量，原因在于申请量能够及时地反映城市创新活动变化。具体而言，本文分别使用发明申请量、实用新型申请量和专利申请总量衡量创新产出，其中发明申请量更接近高质量创新口径，专利申请总量则作为综合创新产出口径。涉及对数形式时，统一使用 $\ln(1+\text{专利量})$ 处理。

核心解释变量为政府产业引导基金规模。基准回归同时考察基金当年设立数量、累计设立数量和累计设立规模。相较于基金数量口径，累计设立规模更直接反映某地截至当年通过引导基金形成的政策性资本供给总量，也更贴近本文关注的政府资本支持强度和资源配置能力，因此本文将其作为正文主线变量，并将基金当年设立数量和累计设立数量作为替代基金口径。

调节变量为地方债务压力。本文采用债券余额合计与公共财政收入之比衡量地方债务压力，该值越高，表示地方偿债压力和财政约束越强。相较于财政收支压力等流量指标，该指标更直接刻画地方政府存量债务负担及其对财政空间、风险承担能力的约束。为缓解同期反向因果和共同冲击的影响，正文主规格优先采用滞后一期债务压力；当期债务压力作为补充口径。财政收支压力更多反映年度预算平衡紧张程度，因此在本文中仅作补充讨论，不作为正文主调节变量。

机制变量围绕三条理论路径设置。第一，耐心资本属性用早期投资事件占比衡量。若投资阶段为种子期或初创期，则记为早期投资，早期投资事件占比等于早期投资事件数除以该市引导基金当年投资总次数。相较于早期投资金额占比，事件占比更直接反映基金是否将投资机会配置到早期项目，也较少受少数大额投资事件影响。第二，社会资本撬动效率用社会资本出资与政府资本出资之比及其变化衡量，其中社会资本出资剔除国资平台、政府产业引导基金和 FOF 后识别。第三，融资约束缓解效应使用地区金融发展水平作为反向代理，正文主口径为存款/GDP，并采用反双曲正弦变换，因为该指标更适合刻画地区金融资源供给能力的存量水平；金融发展增速口径用于补充反映短期改善速度。金融发展水平越高，表示地区金融供给能力越强、融资约束相对越弱，但该指标不能直接等同于企业层面的 SA、KZ 或 WW 融资约束指数。

\begin{table}[htbp]
  \centering
  \caption*{主要变量定义与构造}
  {\fangsong\zihao{-5}
  \begin{tabularx}{0.96\linewidth}{l l X}
    \toprule
    变量类型 & 变量名称 & 构造说明 \\
    \midrule
    创新产出 & 发明申请量 & CNRDS 地方专利数据中地级市—年份发明专利申请数量，用于衡量高质量创新产出。 \\
    创新产出 & 实用新型申请量 & CNRDS 地方专利数据中地级市—年份实用新型专利申请数量，用作补充创新口径。 \\
    创新产出 & 专利申请总量 & 发明、实用新型和外观设计等专利申请数量汇总，用于衡量综合创新产出。 \\
    基金变量 & 基金当年设立数量 & 清科和投中基金目录合并去重后，按注册地和设立年份汇总到城市层面。 \\
    基金变量 & 基金累计设立数量 & 截至当年各地级市累计设立的政府产业引导基金数量。 \\
    基金变量 & 基金累计设立规模 & 截至当年各地级市累计设立的政府产业引导基金规模，正文主解释变量。 \\
    债务压力 & 债务压力 & 债券余额合计除以公共财政收入，数值越高表示地方债务约束越强。 \\
    债务压力 & 滞后一期债务压力 & 债务压力的一期滞后项，正文主调节变量。 \\
    耐心资本 & 早期投资事件占比 & 种子期和初创期投资事件数占该市引导基金当年投资事件总数的比例。 \\
    社会资本 & 社会资本撬动效率 & 非政府社会资本出资与政府资本出资之比，用于衡量政府资本放大效应。 \\
    金融发展 & 存款/GDP & 金融机构存款与地区生产总值之比，用作地区融资约束的反向代理。 \\
    \bottomrule
  \end{tabularx}

  \par\smallskip
  {\songti\zihao{6}注：涉及对数变量时，规模变量统一采用 $\ln(x+1)$ 处理；金融发展主机制采用反双曲正弦变换。}
  }
\end{table}

表1报告主要变量的描述性统计。样本中城市专利申请总量均值为13109.508件，发明申请量均值为4683.297件。基金累计设立规模均值为4.439亿元，但标准差达到16.662亿元，说明政府产业引导基金在城市之间分布差异较大。债务压力均值为233.176，滞后一期债务压力均值为199.112，也显示地方债务约束存在明显地区差异。早期投资事件占比均值为0.340，表明政府引导基金投资事件中约三分之一投向种子期或初创期项目。

\begin{table}[htbp]
  \centering
  \caption{主要变量描述性统计}
  {\fangsong\zihao{-5}
  \resizebox{\linewidth}{!}{%
  \begin{tabular}{l r r r r r}
    \toprule
    变量 & N & 均值 & 标准差 & 最小值 & 最大值 \\
    \midrule
    发明申请量 & 2720 & 4683.297 & 13016.643 & 6.000 & 195388.000 \\
    实用新型申请量 & 2720 & 6420.434 & 13622.132 & 21.000 & 146592.000 \\
    专利申请总量 & 2720 & 13109.508 & 29764.731 & 34.000 & 323014.000 \\
    基金当年设立数量 & 2770 & 2.803 & 7.215 & 0.000 & 99.000 \\
    基金累计设立数量 & 2770 & 14.266 & 39.521 & 0.000 & 458.000 \\
    基金累计设立规模（亿元） & 2629 & 4.439 & 16.662 & 0.000 & 279.266 \\
    债务压力 & 1904 & 233.176 & 484.577 & 0.746 & 4654.410 \\
    滞后一期债务压力 & 1898 & 199.112 & 422.855 & 0.806 & 4129.294 \\
    早期投资事件占比 & 957 & 0.340 & 0.338 & 0.000 & 1.000 \\
    社会资本撬动效率 & 986 & 0.789 & 2.015 & 0.000 & 34.050 \\
    金融发展水平 & 2500 & 1.638 & 0.772 & 0.000 & 20.100 \\
    经济发展水平 & 2098 & 5.453 & 0.864 & 3.081 & 8.590 \\
    财政科技支出 & 1971 & 6.343 & 1.418 & 2.757 & 11.060 \\
    人口规模 & 2450 & 5.980 & 0.748 & 3.252 & 11.975 \\
    第二产业规模 & 2098 & 4.554 & 0.886 & 2.272 & 7.238 \\
    外资利用水平 & 1681 & 10.184 & 2.082 & 1.386 & 14.776 \\
    \bottomrule
  \end{tabular}%
  }

  \par\smallskip
  {\songti\zihao{6}注：基金累计设立规模按亿元报告。经济发展水平、财政科技支出、人口规模、第二产业规模和外资利用水平为 $\ln(x+1)$ 口径。社会资本撬动效率报告有限值观测，剔除因政府出资为0形成的非有限值。}
  }
\end{table}

控制变量沿用现有回归设定，主要包括经济发展水平、财政科技支出、人口规模、产业结构和外资利用水平。经济发展水平用地区生产总值对数衡量，财政科技支出用于控制地方财政对科技创新的直接支持强度，常住人口对数用于控制人口规模和市场规模，第二产业增加值对数用于控制城市产业基础，实际利用外资额对数用于控制对外开放程度。上述控制变量均按 $\ln(x+1)$ 处理，以缓解规模变量右偏分布并保留取值为零的观测。

\subsection{模型设定}

为检验政府产业引导基金是否促进城市创新，本文首先估计如下双向固定效应模型：

\begin{equation}
Innovation_{it}=\alpha+\beta Fund_{it}+\gamma Controls_{it}+\mu_i+\lambda_t+\varepsilon_{it}.
\end{equation}
\ERJWhere{$Innovation_{it}$ 表示城市 $i$ 在年份 $t$ 的创新产出，$Fund_{it}$ 表示政府产业引导基金变量，$Controls_{it}$ 为城市层面控制变量，$\mu_i$ 和 $\lambda_t$ 分别表示城市固定效应和年份固定效应。若 $\beta$ 显著为正，则说明政府产业引导基金扩张与城市创新产出提升相一致。}

在此基础上，本文引入债务压力与基金规模的交互项，考察地方债务压力是否削弱基金的创新促进效应：

\begin{equation}
Innovation_{it}=\alpha+\beta_1 Fund_{it}+\beta_2 Debt_{i,t-1}
+\beta_3(Fund_{it}\times Debt_{i,t-1})+\gamma Controls_{it}+\mu_i+\lambda_t+\varepsilon_{it}.
\end{equation}
\ERJWhere{$Debt_{i,t-1}$ 表示滞后一期债务压力，$Fund_{it}\times Debt_{i,t-1}$ 为核心交互项。本文关注 $\beta_3$ 的符号和显著性。若 $\beta_3<0$，说明债务压力越高，同等规模政府产业引导基金对创新产出的边际促进作用越弱。}

机制检验围绕耐心资本属性、社会资本撬动效率和融资约束缓解效应展开。耐心资本和金融发展机制主要采用“债务压力削弱基金影响机制变量，机制变量进一步影响创新产出”的路径检验：

\begin{equation}
M_{it}=\alpha+\theta_1(Fund_{it}\times Debt_{i,t-1})+\gamma Controls_{it}+\mu_i+\lambda_t+\nu_{it},
\end{equation}

\begin{equation}
Innovation_{i,t+k}=\alpha+\rho_1(Fund_{it}\times Debt_{i,t-1})+\rho_2 M_{it}
+\gamma Controls_{it}+\mu_i+\lambda_t+\eta_{it}.
\end{equation}
\ERJWhere{$M_{it}$ 表示机制变量，$k$ 表示创新产出的滞后实现期。耐心资本机制中 $M_{it}$ 为早期投资事件占比，理论预期为 $\theta_1<0$ 且 $\rho_2>0$；金融发展机制中 $M_{it}$ 为金融发展水平，理论预期同样为 $\theta_1<0$ 且 $\rho_2>0$。}

社会资本撬动效率机制不写作普通中介，而写作调节式机制或承接机制。相应模型在结果方程中进一步加入基金规模与社会资本撬动效率的交互项：

\begin{equation}
Innovation_{it}=\alpha+\delta_1(Fund_{it}\times Debt_{i,t-1})
+\delta_2(Fund_{it}\times SocCap_{it})+\gamma Controls_{it}+\mu_i+\lambda_t+\xi_{it}.
\end{equation}
\ERJWhere{$SocCap_{it}$ 表示社会资本撬动效率或其变化。若债务压力降低撬动效率，同时 $\delta_2>0$，且加入机制项后 $\delta_1$ 的绝对值下降，则说明债务压力对基金创新效应的抑制部分通过社会资本撬动效率这一承接渠道体现。}

以上模型均采用城市固定效应和年份固定效应，并按城市维度聚类稳健标准误。本文的识别重点不是宣称完全解决所有内生性问题，而是在固定城市不随时间变化的特征、年份共同冲击以及主要城市层面控制变量后，比较同一城市在不同年份基金规模、债务压力和创新产出的共同变化关系。后续滞后债务压力和滞后创新响应的检验，用于进一步缓解反向因果和同期冲击的担忧。


% 本文件由 main.tex 输入，作为“实证结果分析”章节，不单独编译。

\section{实证结果分析}

\subsection{基准回归：政府产业引导基金与城市创新}

基准回归首先检验政府产业引导基金是否能够提升城市创新产出。表2报告了核心结果。可以看到，在控制城市固定效应、年份固定效应以及经济发展水平、财政科技支出、人口规模、产业结构和外资利用水平后，政府产业引导基金设立数量和累计设立规模的系数总体显著为正。其中，累计设立规模对发明申请量、实用新型申请量和专利申请总量均显著为正，说明基金规模扩张与城市专利申请产出提高存在稳定联系。

\begin{table}[htbp]
  \centering
  \caption{政府产业引导基金与城市创新：基准回归}
  {\fangsong\zihao{-5}
  \begin{tabularx}{0.96\linewidth}{l l c c c}
    \toprule
    因变量 & 基金口径 & 系数 & p值 & N \\
    \midrule
    发明申请量 & 累计设立规模 & 0.0300*** & $<0.001$ & 1083 \\
    实用新型申请量 & 累计设立规模 & 0.0405*** & 0.0031 & 1083 \\
    专利申请总量 & 累计设立规模 & 0.0800** & 0.0174 & 1083 \\
    发明申请量 & 当年设立数量 & 297.1917*** & $<0.001$ & 1149 \\
    专利申请总量 & 当年设立数量 & 849.5984*** & 0.0018 & 1149 \\
    发明申请量 & 累计设立数量 & 106.7968*** & $<0.001$ & 1149 \\
    \midrule
    控制变量 & \multicolumn{4}{c}{是} \\
    城市固定效应 & \multicolumn{4}{c}{是} \\
    年份固定效应 & \multicolumn{4}{c}{是} \\
    \bottomrule
  \end{tabularx}

  \par\smallskip
  {\songti\zihao{6}注：所有模型均按城市聚类稳健标准误。*、**、*** 分别表示在10\%、5\%、1\% 水平上显著。}
  }
\end{table}

% TODO：结果 CSV 解锁后，可按 README 要求为表2和表3补入各规格 $R^2$。

从经济含义看，基金累计设立规模对发明申请量和专利申请总量均显著为正，表明政府产业引导基金的作用并不限于增加一般性专利申请，也与高质量创新口径的改善相一致。当年设立数量和累计设立数量同样表现出显著正向结果，说明无论从基金数量还是基金规模观察，引导基金扩张都与城市创新产出提升相一致。由此，假说 H1 获得较强支持。考虑到累计设立规模在后续债务压力调节效应和机制检验中更为连贯，下文将其作为正文主解释变量。

\subsection{债务压力的调节效应}

在确认基金总体上具有创新促进效应后，本文进一步检验地方债务压力是否削弱这一作用。表3报告基金累计设立规模与滞后一期债务压力交互项的估计结果。交互项在专利申请总量、发明申请量和实用新型申请量三个口径下均显著为负，说明债务压力较高的城市中，同等规模政府产业引导基金所能带来的创新促进作用明显下降。

\begin{table}[htbp]
  \centering
  \caption{债务压力对基金创新效应的调节作用}
  {\fangsong\zihao{-5}
  \begin{tabularx}{0.92\linewidth}{l c c c}
    \toprule
    因变量 & 交互项系数 & p值 & N \\
    \midrule
    专利申请总量 & -0.0000609*** & $<0.001$ & 914 \\
    发明申请量 & -0.0000162*** & $<0.001$ & 914 \\
    实用新型申请量 & -0.0000292*** & $<0.001$ & 914 \\
    \midrule
    控制变量 & \multicolumn{3}{c}{是} \\
    城市固定效应 & \multicolumn{3}{c}{是} \\
    年份固定效应 & \multicolumn{3}{c}{是} \\
    \bottomrule
  \end{tabularx}

  \par\smallskip
  {\songti\zihao{6}注：交互项为“基金累计设立规模$\times$滞后一期债务压力”。所有模型均按城市聚类稳健标准误。同期债务压力口径下交互项同样显著为负。财政压力交互项在现有结果中方向与理论主线不一致，因此不纳入正文主结论。}
  }
\end{table}

该结果支持假说 H2。其含义是，政府产业引导基金并不是在任何财政环境下都能发挥同等作用。当地方债务压力上升时，地方政府的财政空间、风险承担能力和政策耐心可能下降，基金更容易受到短期绩效压力和偿债约束影响，从而削弱其对创新活动的边际促进作用。由于滞后一期债务压力更符合“债务约束先于基金绩效变化”的时间顺序，本文将表3作为调节效应的主规格。

\subsection{稳健性检验：替换变量口径}

为检验主结论是否依赖单一变量口径，本文从创新产出、基金规模和债务压力三个方面进行替换。首先，创新产出不仅使用专利申请总量，也使用发明申请量和实用新型申请量。基准回归中，基金累计设立规模对三类创新产出的系数均显著为正；调节效应中，基金累计设立规模与债务压力的交互项在三类创新产出下均显著为负。其次，基金变量除累计设立规模外，还使用当年设立数量和累计设立数量。结果显示，基金数量口径同样与城市创新产出显著正相关。最后，债务压力除滞后一期口径外，也使用当期债务压力和缩尾债务压力作为补充，核心交互项仍保持负向显著。

机制检验中的变量替换也提供了补充支持。耐心资本机制中，早期投资事件占比采用原始比例和1/99缩尾比例后均得到一致结果；金融发展机制中，存款/GDP 的反双曲正弦变换是正文主口径，对数变换口径方向相同但显著性略弱。综合来看，基准结论和债务压力调节结论并非由单一创新变量、单一基金变量或单一债务压力口径驱动。

\subsection{内生性处理：滞后变量检验}

本文进一步采用滞后变量缓解反向因果和同期冲击的担忧。债务压力调节效应的正文主规格使用滞后一期债务压力，而不是仅使用当期债务压力。表3显示，基金累计设立规模与滞后一期债务压力的交互项在三类创新产出下均显著为负，说明即便让债务压力先于创新结果进入模型，债务压力削弱基金创新效应的结论仍然成立。

机制检验也体现了创新产出的滞后响应。耐心资本机制中，早期投资事件占比主要影响未来两年专利申请总量；这与早期项目从投资到形成专利产出的时间滞后相一致。需要强调的是，滞后变量并不能完全消除所有内生性问题，尤其不能排除未观测的时变政策冲击。但在城市固定效应、年份固定效应和控制变量基础上，滞后债务压力与滞后创新响应结果增强了本文主结论的可信度。

\subsection{机制检验一：耐心资本属性}

第一条机制关注政府产业引导基金的耐心资本属性。理论上，债务压力越高，地方政府越可能要求基金更快见效、更低风险地配置资金，从而削弱其“投早、投小、投长期”的功能。本文使用早期投资事件占比代理耐心资本属性。为避免比例变量分母过小导致机械波动，机制检验将样本限定为基金投资事件数不少于2的城市—年份观测。该处理是识别比例机制的必要步骤，因为当分母为1时，单次投资会使早期投资占比机械地落在0或1，容易把偶然事件误判为结构性投资偏好。

\begin{table}[htbp]
  \centering
  \caption{耐心资本属性机制：早期投资事件占比}
  {\fangsong\zihao{-5}
  \resizebox{\linewidth}{!}{%
  \begin{tabular}{l l l c c c c c}
    \toprule
    创新结果 & 比例处理 & 债务处理 & 总效应 & 路径一 & 路径二 & $R^2$ & N \\
    \midrule
    两年后专利申请总量 & 原始比例 & 原始债务压力 & -0.085** & -0.143** & 0.092** & 0.680 & 266 \\
    两年后专利申请总量 & 比例1/99缩尾 & 原始债务压力 & -0.085** & -0.143** & 0.092** & 0.680 & 266 \\
    两年后专利申请总量 & 原始比例 & 债务压力1/99缩尾 & -0.075** & -0.128** & 0.091** & 0.680 & 266 \\
    两年后专利申请总量 & 比例1/99缩尾 & 债务压力1/99缩尾 & -0.075** & -0.128** & 0.091** & 0.680 & 266 \\
    \midrule
    控制变量 & \multicolumn{7}{c}{是} \\
    城市固定效应 & \multicolumn{7}{c}{是} \\
    年份固定效应 & \multicolumn{7}{c}{是} \\
    \bottomrule
  \end{tabular}%
  }

  \par\smallskip
  {\songti\zihao{6}注：总效应表示“基金规模$\times$债务压力”对未来创新的影响，路径一表示该交互项对早期投资事件占比的影响，路径二表示早期投资事件占比对未来创新的影响。表中结果来自基金投资事件数不少于2的分母稳定化样本。}
  }
\end{table}

表4显示，基金规模与债务压力的交互项显著降低早期投资事件占比，而早期投资事件占比对未来两年专利申请总量具有显著正向影响。与此同时，总效应显著为负。这说明债务压力可能通过削弱引导基金投向早期项目的耐心资本属性，进一步抑制其创新促进作用。需要谨慎的是，该组结果依赖分母稳定化样本，且早期投资金额占比口径并不稳定。因此，本文将其解释为具有透明样本边界的机制证据，而不将其写作无条件结论。

\subsection{机制检验二：社会资本撬动效率}

第二条机制关注社会资本撬动效率。政府产业引导基金的重要功能不仅在于直接提供财政资金，也在于通过政府背书和资本放大效应带动社会资本进入创新领域。债务压力上升时，地方政府可能降低基金运行稳定性、市场化程度和风险承担能力，从而削弱社会资本跟投和撬动效率。现有结果并不支持“社会资本撬动效率是普遍普通中介”的表述，而更支持一种调节式或承接机制：债务压力削弱撬动效率及其改善，而较高撬动效率会强化基金规模对创新的边际促进作用。

\begin{table}[htbp]
  \centering
  \caption{社会资本撬动效率机制：低金融发展地区}
  {\fangsong\zihao{-5}
  \resizebox{\linewidth}{!}{%
  \begin{tabular}{l l c c c c c}
    \toprule
    因变量 & 债务处理 & 基准交互项 & 机制方程 & 结果方程 & 加机制后交互项 & N \\
    \midrule
    实用新型申请量对数 & 滞后债务压力 & -4.03E-08*** & -8.95E-08** & 6.31E-07** & -3.68E-08*** & 705 \\
    实用新型申请量对数 & 滞后债务压力缩尾 & -4.03E-08*** & -8.95E-08** & 6.31E-07** & -3.68E-08*** & 705 \\
    下一期专利申请总量 & 滞后债务压力 & -8.20E-04** & -8.95E-08** & 0.02418** & -6.61E-04*** & 705 \\
    下一期专利申请总量 & 滞后债务压力缩尾 & -8.20E-04** & -8.95E-08** & 0.02418** & -6.61E-04*** & 705 \\
    \bottomrule
  \end{tabular}%
  }

  \par\smallskip
  {\songti\zihao{6}注：表中结果来自低金融发展、剔除2020年的样本，规格为未加控制变量的5\%全链条结果。机制变量为社会资本撬动效率年度变化，缺失补0后缩尾。加入机制项后，原债务调节项绝对值下降。}
  }
\end{table}

表5显示，在金融发展水平较低的城市中，债务压力会削弱政府引导基金撬动社会资本的效率或效率改善；而社会资本撬动效率越高，基金规模对创新产出的边际促进作用越强。加入社会资本撬动效率相关机制项后，原基金规模与债务压力交互项的绝对值下降，说明债务压力的抑制作用部分通过社会资本撬动效率这一承接渠道体现。由于最强证据集中在低金融发展地区，且5\%全链条结果来自未加控制规格，本文对该机制采用保守表述，将其限定为特定金融环境下更明显的调节式机制。

\subsection{机制检验三：融资约束缓解效应}

第三条机制检验政府产业引导基金能否通过改善地区金融发展环境来缓解融资约束。本文使用地区金融发展水平作为融资约束的反向代理，而不将其等同于企业层面的融资约束指数。在所有可用链条中，最适合正文报告的是存款/GDP 口径金融发展水平的反双曲正弦变换结果。

\begin{table}[htbp]
  \centering
  \caption{融资约束缓解机制：金融发展水平}
  {\fangsong\zihao{-5}
  \resizebox{\linewidth}{!}{%
  \begin{tabular}{l l c c c c c}
    \toprule
    机制变量口径 & 创新结果 & 交互项$\rightarrow$机制 & p值 & 机制$\rightarrow$创新 & p值 & N \\
    \midrule
    存款/GDP水平（asinh变换） & 专利申请总量对数 & -0.000000033 & 0.0147 & 0.4039 & 0.0497 & 944 \\
    贷款/GDP增速 & 实用新型申请量 & -0.001032 & 0.0228 & 10986.6163 & 0.0336 & 782 \\
    \bottomrule
  \end{tabular}%
  }

  \par\smallskip
  {\songti\zihao{6}注：第一行为正文主链条，表示基金累计设立规模与滞后债务压力交互项削弱金融发展水平，进而不利于创新产出。第二行为补充稳健性口径，其含义是金融发展改善速度。所有结果均纳入控制变量、城市固定效应和年份固定效应，并按城市聚类稳健标准误。}
  }
\end{table}

表6表明，当以存款/GDP 衡量金融发展水平并进行反双曲正弦变换时，基金累计设立规模与滞后债务压力的交互项对金融发展水平的影响显著为负，而金融发展水平对总专利申请对数显著为正。这说明债务压力会削弱政府产业引导基金改善地区金融环境、缓解融资约束的能力，并由此抑制城市创新产出。需要注意的是，主链条第二步的显著性接近5\%临界值，因此这一机制虽然与理论预期一致，但仍应以谨慎口径表述。

\subsection{本章小结}

综合来看，政府产业引导基金能够显著提升城市创新产出，但这一促进作用受到地方债务压力的制约。债务压力越高，基金累计设立规模对专利申请的边际促进效应越弱。进一步的机制检验表明，债务压力可能通过三条路径削弱引导基金的创新效应：一是降低基金投向早期项目的比例，削弱其耐心资本属性；二是在金融发展水平较低地区降低基金撬动社会资本的效率；三是削弱基金改善地区金融发展、缓解融资约束的能力。

上述结果也存在清晰边界。耐心资本机制依赖于基金投资事件数不少于2的分母稳定化样本；社会资本撬动效率机制主要出现在低金融发展地区，属于情境依赖较强的承接机制；金融发展机制虽然在理论链条上较为完整，但关键显著性处于5\%边界附近。因此，本文更适合将机制结果理解为与主结论方向一致、但具有样本条件和识别边界的补充证据。总体而言，政府产业引导基金的创新促进作用并非无条件成立，其政策绩效在很大程度上取决于地方财政债务约束与地区资本市场环境。


\section{结论与进一步讨论}

这里填写整篇论文的总体结论与政策含义。当前实证章节末尾已经给出本节小结，本节可在引言与理论分析补齐之后，再统一整合为全文结论。

\section*{参考文献}

{\songti\zihao{6}
\ERJRef{蔡庆丰、刘昊、舒少文，2024：《政府产业引导基金与域内企业创新：引导效应还是挤出效应？》，《财贸经济》第3期。}
\ERJRef{刘贯春、程飞阳、姚守宇、张军，2022：《地方政府债务治理与企业投融资期限错配改善》，《管理世界》第11期。}
\ERJRef{刘玉斌、赵天宇、郭树龙，2023：《战略性新兴产业创业投资引导基金能促进企业创新吗？》，《产业经济研究》第1期。}
\ERJRef{施国平、党兴华、董建卫，2016：《引导基金能引导创投机构投向早期和高科技企业吗？——基于双重差分模型的实证评估》，《科学学研究》第6期。}
\ERJRef{王平、徐肇仪，2023：《数字普惠金融、融资约束与共同富裕》，《会计之友》第10期。}
\ERJRef{吴超鹏、严泽浩，2023：《政府基金引导与企业核心技术突破：机制与效应》，《经济研究》第6期。}
\ERJRef{徐明，2022：《政府风险投资、代理问题与企业创新——来自政府引导基金介入的证据》，《南开经济研究》第2期。}
\ERJRef{徐彦坤，2020：《地方政府债务如何影响企业投融资行为？》，《中南财经政法大学学报》第2期。}
\ERJRef{张杰、任元明，2025：《探寻以耐心资本为内核的中国特色现代金融体系发展路径与改革突破口》，《天津社会科学》第2期。}
\ERJRef{张壹帆、陆岷峰，2025：《新质生产力视角下政府引导基金“耐心度”培育：评价体系与利益平衡》，《证券市场导报》第1期。}
\ERJRef{张跃文、李慧、徐枫，2025：《财政压力、“合作型”政府干预与企业投资》，《财经论丛》第9期。}
\ERJRef{周伟、陆佳玮，2022：《财政压力、地方政府行为与企业创新》，出版信息待核。}
\ERJRef{Cumming, D. J., 2007, “Government Policy towards Entrepreneurial Finance: Innovation Investment Funds”, Journal of Business Venturing, Vol. 22, No. 2, pp. 193-235.}
\ERJRef{Cumming, D. J., and S. Johan, 2007, “Pre-Seed Government Venture Capital Funds”, Journal of International Entrepreneurship, forthcoming.}
}

\vspace{2\baselineskip}
\begin{flushright}
（责任编辑：\underline{\hspace{4em}}）（校对：\underline{\hspace{4em}}）
\end{flushright}

\end{document}
